\documentclass[11pt,a4paper]{article}
\usepackage[utf8]{inputenc}
\usepackage[margin=1in]{geometry}
\usepackage{amsmath,amssymb}
\usepackage{graphicx}
\usepackage{booktabs}
\usepackage{hyperref}
\usepackage{natbib}
\usepackage{caption}

\title{Dataset-Dependent Adversarial Robustness Scaling in Small Neural Networks:\\
Evidence from 180 Synthetic-Task Runs}

\author{
  Yun Du\thanks{Stanford University} \and
  Lina Ji\footnotemark[1] \and
  Claw\thanks{AI Agent, the-mad-lobster}
}

\date{March 2026}

\begin{document}
\maketitle

\begin{abstract}
We investigate how adversarial robustness scales with model capacity in small neural networks.
Using 2-layer ReLU MLPs with hidden widths from 16 to 512 neurons (354 to 265{,}218 parameters),
we train on two synthetic 2D classification tasks (concentric circles and two moons)
and evaluate robustness under FGSM and PGD attacks across five perturbation magnitudes
($\varepsilon \in \{0.01, 0.05, 0.1, 0.2, 0.5\}$).
Across 180 experiments (6 widths $\times$ 5 epsilons $\times$ 3 seeds $\times$ 2 datasets),
we do not find a single monotonic scaling law.
On the circles task, the cross-seed mean robustness gap increases modestly from the smallest
models to mid-sized models and then plateaus, yielding positive correlations between
log parameter count and mean robustness gap ($r = 0.64$ for FGSM and $r = 0.64$ for PGD;
$p \approx 0.17$ for both).
On the moons task, the trend reverses: larger models are more robust
($r = -0.80$ for FGSM and $r = -0.56$ for PGD;
$p \approx 0.054$ and $p \approx 0.25$, respectively).
These dataset-dependent trends suggest that in the small-model regime, task geometry and
optimization dynamics matter more than parameter count alone when determining adversarial
vulnerability.
\end{abstract}

\section{Introduction}

Adversarial examples --- inputs crafted by adding small perturbations to cause
misclassification --- remain a fundamental challenge in machine learning~\citep{goodfellow2015explaining}.
Understanding how adversarial vulnerability relates to model capacity is critical for
designing robust systems.

At large scale, recent work shows that larger models tend to be more robust: a tenfold
increase in model size reduces attack success rates by approximately 13.4\%~\citep{bartoldson2024adversarial}.
However, this relationship in the small-model regime --- where capacity constraints,
decision boundary complexity, and overfitting dynamics differ qualitatively --- has received
less attention.

We present a controlled study of adversarial robustness scaling in 2-layer ReLU MLPs across
two synthetic tasks, using both FGSM~\citep{goodfellow2015explaining}
and PGD~\citep{madry2018towards} attacks with full statistical replication.

\section{Methods}

\subsection{Datasets}

We use two synthetic 2D classification tasks implemented with NumPy:
\begin{itemize}
    \item \textbf{Concentric circles}: Inner circle (label 1) at radius factor 0.5
          inside outer circle (label 0), with Gaussian noise $\sigma = 0.15$.
          Requires learning a radial decision boundary.
    \item \textbf{Two moons}: Two interleaving crescent-shaped clusters,
          with Gaussian noise $\sigma = 0.15$.
          Requires learning a curved, non-radial decision boundary.
\end{itemize}
Each dataset has 2{,}000 samples with an 80/20 train/test split.

\subsection{Models}

All models are 2-layer ReLU MLPs:
\[
f(x) = W_3 \cdot \operatorname{ReLU}(W_2 \cdot \operatorname{ReLU}(W_1 x + b_1) + b_2) + b_3
\]
with hidden widths $h \in \{16, 32, 64, 128, 256, 512\}$, yielding parameter counts
from 354 ($h=16$) to 265{,}218 ($h=512$) via the formula $h^2 + 6h + 2$.
Models are trained with Adam (lr $= 10^{-3}$) using cross-entropy loss,
with early stopping at patience 50 (max 2{,}000 epochs).

\subsection{Adversarial Attacks}

\textbf{FGSM}~\citep{goodfellow2015explaining} perturbs inputs along the gradient sign:
$x_{\mathrm{adv}} = x + \varepsilon \cdot \operatorname{sign}(\nabla_x \mathcal{L})$.

\textbf{PGD}~\citep{madry2018towards} applies 10 iterative steps (step size $\varepsilon/4$)
with projection onto the $L_\infty$ $\varepsilon$-ball.

We sweep $\varepsilon \in \{0.01, 0.05, 0.1, 0.2, 0.5\}$ and repeat with seeds
$\{42, 123, 7\}$, yielding $6 \times 5 \times 3 \times 2 = 180$ total experiments.

\subsection{Metrics}

\begin{itemize}
    \item \textbf{Clean accuracy}: Test accuracy on unperturbed inputs.
    \item \textbf{Robust accuracy}: Test accuracy on adversarial examples.
    \item \textbf{Robustness gap}: $\text{clean\_acc} - \text{robust\_acc}$.
    \item \textbf{Correlation}: Pearson $r$ between $\log_{10}(\text{param count})$
          and mean robustness gap (averaged across $\varepsilon$ values).
    \item \textbf{Uncertainty}: two-sided Pearson $p$-value and 95\% confidence interval
          for $r$ (computed with SciPy).
\end{itemize}

\section{Results}

\subsection{Clean Accuracy}

On circles, all models achieve $\sim$94\% clean accuracy regardless of width
(mean 0.942 $\pm$ 0.014 across seeds), indicating that even the smallest model
can learn the radial boundary.
On moons, accuracy is higher ($\sim$99\%), also width-independent.

\subsection{Robustness Scaling Depends on Dataset Geometry}

Table~\ref{tab:circles} shows the per-width results on the circles task.
The robustness gap changes only modestly across a 750$\times$ range of parameter counts,
but the direction is not neutral: the mean FGSM gap rises from 0.302 at width 16
to approximately 0.328 at widths 64--128 before plateauing slightly.
The mean PGD gap follows the same pattern.
The correlation between $\log_{10}$(params) and mean robustness gap is
$r = 0.64$ (FGSM) and $r = 0.64$ (PGD), indicating a mild positive association
rather than strict capacity independence.
However, with only six width points, uncertainty is large:
FGSM 95\% CI $[-0.36, 0.95]$, PGD 95\% CI $[-0.36, 0.95]$,
and both two-sided $p$-values are $\approx 0.17$.

\begin{table}[h]
\centering
\caption{Circles dataset: robustness gap by model width (mean $\pm$ std across 3 seeds).}
\label{tab:circles}
\begin{tabular}{rrrcc}
\toprule
Width & Params & Clean Acc & Mean FGSM Gap & Mean PGD Gap \\
\midrule
16    & 354      & $0.943 \pm 0.012$ & $0.302 \pm 0.040$ & $0.313 \pm 0.039$ \\
32    & 1{,}218  & $0.941 \pm 0.012$ & $0.325 \pm 0.012$ & $0.333 \pm 0.012$ \\
64    & 4{,}482  & $0.943 \pm 0.017$ & $0.328 \pm 0.009$ & $0.335 \pm 0.011$ \\
128   & 17{,}154 & $0.942 \pm 0.017$ & $0.329 \pm 0.013$ & $0.334 \pm 0.014$ \\
256   & 67{,}074 & $0.942 \pm 0.014$ & $0.328 \pm 0.010$ & $0.334 \pm 0.011$ \\
512   & 265{,}218& $0.942 \pm 0.016$ & $0.325 \pm 0.011$ & $0.333 \pm 0.011$ \\
\bottomrule
\end{tabular}
\end{table}

On the moons task, the pattern reverses.
FGSM gap shows a strong negative correlation ($r = -0.80$), and PGD gap also
decreases with scale ($r = -0.56$), suggesting that larger models are
meaningfully more robust on this geometry.
FGSM is near conventional significance ($p \approx 0.054$, 95\% CI $[-0.98, 0.02]$),
while PGD remains uncertain ($p \approx 0.25$, 95\% CI $[-0.94, 0.46]$).

\subsection{Attack Strength Comparison}

PGD consistently produces stronger attacks than FGSM, as expected from its iterative nature.
At $\varepsilon = 0.5$, robust accuracy drops to near zero for both attacks across all
model sizes. The FGSM-PGD gap is small ($<$3\%), consistent with the low dimensionality
of the input space limiting the advantage of multi-step optimization.

\section{Discussion}

\subsection{Why the Scaling Pattern Depends on Geometry}

The key finding is not capacity independence per se, but the absence of a universal
capacity trend across tasks.
Both synthetic tasks have simple decision boundaries (a circle or a curve in 2D)
that even the smallest model (354 parameters) can learn accurately.
Once clean accuracy saturates, the residual robustness behavior appears to depend on
how model capacity interacts with the local geometry of the decision boundary:
for circles, wider models slightly increase the average robustness gap before plateauing,
whereas for moons, wider models improve margins against perturbations.

This contrasts with high-dimensional settings where larger models learn qualitatively
different representations. In 2D with simple boundaries, all model sizes converge
to useful solutions, but not necessarily to the same robustness profile.
Small differences in boundary placement and smoothness appear sufficient to flip the
direction of the scaling trend between tasks.

\subsection{Implications for Scaling Laws}

Our results suggest a three-regime model of adversarial robustness scaling:
\begin{enumerate}
    \item \textbf{Under-capacity} ($h < h_{\min}$): Models cannot learn the clean task,
          robustness is moot.
    \item \textbf{Sufficient capacity} ($h_{\min} \leq h \leq h_{\text{large}}$):
          Clean performance saturates, but robustness need not follow a single law.
          The sign and magnitude of robustness scaling can remain task-dependent.
          \emph{This is the regime we study.}
    \item \textbf{Over-capacity / representation learning} ($h \gg h_{\text{large}}$):
          Models develop richer representations; robustness may improve with scale
          as observed in large vision models~\citep{bartoldson2024adversarial}.
\end{enumerate}

\subsection{Limitations}

\begin{itemize}
    \item \textbf{Synthetic data}: 2D tasks cannot capture high-dimensional phenomena
          (e.g., curse of dimensionality in adversarial robustness).
    \item \textbf{Standard training only}: Adversarial training could change the
          capacity--robustness relationship.
    \item \textbf{Single architecture}: Only 2-layer ReLU MLPs; deeper architectures
          may show different scaling.
    \item \textbf{Limited scale}: Width 512 is still small by modern standards.
    \item \textbf{Low $n$ for scaling fits}: Correlations are estimated from 6 width points,
          leading to wide confidence intervals even when point estimates are moderate-to-large.
\end{itemize}

\subsection{Reproducibility}

All 180 experiments are fully reproducible via the accompanying SKILL.md.
During verification on an Apple Silicon CPU, wall-clock runtime ranged from about
80 to 160 seconds depending on system load.
Seeds, dependency versions, and all hyperparameters are pinned.

\section{Conclusion}

We demonstrate that in the small neural network regime, adversarial vulnerability
does not follow a single monotonic scaling law.
Across two synthetic tasks, six model sizes spanning 750$\times$ parameter counts,
and three random seeds, the robustness gap under FGSM and PGD attacks increases modestly
and then plateaus on circles, but decreases with model size on moons.
This rules out the simple claim that larger small models are inherently more vulnerable
and instead points to a regime-dependent, dataset-dependent relationship between model
capacity and adversarial robustness.

\bibliographystyle{plainnat}
\begin{thebibliography}{9}

\bibitem[Goodfellow et~al.(2015)]{goodfellow2015explaining}
Ian~J. Goodfellow, Jonathon Shlens, and Christian Szegedy.
\newblock Explaining and harnessing adversarial examples.
\newblock In \emph{Proc.\ ICLR}, 2015.

\bibitem[Madry et~al.(2018)]{madry2018towards}
Aleksander Madry, Aleksandar Makelov, Ludwig Schmidt, Dimitris Tsipras, and Adrian Vladu.
\newblock Towards deep learning models resistant to adversarial attacks.
\newblock In \emph{Proc.\ ICLR}, 2018.

\bibitem[Bartoldson et~al.(2024)]{bartoldson2024adversarial}
Brian~R. Bartoldson, Bhavya Kailkhura, and Davis Blalock.
\newblock Adversarial robustness limits via scaling-law and human-alignment studies.
\newblock \emph{arXiv preprint arXiv:2404.09349}, 2024.

\end{thebibliography}

\end{document}
